\section{Observable compatibility and a geometric application}
\label{sec:v84-observed-geometry}
The common-lift assumption can be replaced by identities of the
reconstructed objects, provided one works on an observed invariant
domain. We first specify what is to be checked and what domain is
actually obtained. An identity of phase maps alone is insufficient:
the absolute primitives are part of every identity below.

\subsection{Graph consistency and domains of words}
For each indicated count $n$, recover all absolute sheets and form
\[
 \mathcal G_n=\bigcup_{\text{gates, sheets}}
 \{((x,-S_x(x,y)),(y,S_y(x,y)),S(x,y))\}.
\]
Regularity makes each piece a local exact lift. On an open subset of
the union of its source projections, check that all graph pieces with
the same source have the same target \emph{and the same primitive}.
Check likewise that the target projection has only one source. These
are equalities and injectivity tests on the recovered graphs. Where
they hold, the pieces glue to an exact diffeomorphism between the
observed source and target domains; no semantic component matching is
supplied. The tests may fail for inconsistent data, in which case no
common deterministic graph is asserted on that region.

For a finite word with factors $a_1,\ldots,a_N$, in their order of
application, write the observed partial lifts as $(F_i,A_i)$ and
let $D_i$ be their observed source domains. For an inverse factor the
source is the corresponding observed range and the primitive is
$-A\circ F^{-1}$. Define recursively
\begin{equation}\label{eq:v84-word-domain}
 \begin{split}
 \Omega_0&=X,\qquad Z_0=\id,\quad B_0=0,\\
 \Omega_i&=\{z\in\Omega_{i-1}:Z_{i-1}(z)\in D_i\},\\
 Z_i&=F_i\circ Z_{i-1},\qquad
 B_i=B_{i-1}+A_i\circ Z_{i-1}\quad\hbox{on }\Omega_i.
 \end{split}
\end{equation}
Each $\Omega_i$ is open. Induction on the number of factors proves
that $\Omega_N$ is exactly the domain on which this word can be
evaluated using these observed pieces. The formula keeps every
primitive constant and resolves the order of negative factors.
For example, $g^3(g^2)^{-1}$ first uses the observed inverse of $g^2$,
then the observed $g^3$ at its output. It is not enough for both maps
to be separately known somewhere on the original target.
This is a maximal \emph{word-evaluation} domain, not a claim of
maximal recovery by all possible geometric arguments.

\begin{theorem}[Certification of a common return lift]
\label{thm:v84-compatible}
Let $r,s$ be coprime positive integers. Suppose graph consistency and
source--target tests on the recovered count-$r$ and count-$s$ sheets
give exact lifts $R,S$ that are diffeomorphisms of an observed open
domain $\Omega$ onto itself. There exists an exact lift $g$ on $\Omega$
with $g^r=R$ and $g^s=S$ if and only if
\begin{equation}\label{eq:v84-compatibility}
                         RS=SR,\qquad R^s=S^r
\end{equation}
as \emph{full exact lifts}. When these conditions hold, that lift is
unique and, for any $ur+vs=1$,
\[
                              g=R^uS^v.
\]
Thus the common-lift premise is testable from the reconstructed graphs
and primitives on $\Omega$. The counts, canonical coordinates and
apparatus class remain supplied information; a section or unobserved
phase region is not inferred by this test.
\end{theorem}
\begin{proof}
Powers of a common lift satisfy \eqref{eq:v84-compatibility}. Conversely,
exact lifts form a group under \eqref{eq:v83-lifts} on the invariant
domain. Commutativity and $R^s=S^r$ give
\[
 (R^uS^v)^r=R^{ur}(S^r)^v=R^{ur+sv}=R,
 \qquad
 (R^uS^v)^s=(R^s)^uS^{vs}=S^{ru+vs}=S.
\]
These equalities hold in the lift group, so include the primitives.
If $h$ is any other common lift, then
$h=h^{ur+vs}=R^uS^v$, proving uniqueness and independence of the
Bezout pair. All compositions are defined because both lifts preserve
$\Omega$. On noninvariant observed domains one uses
\eqref{eq:v84-word-domain} instead and makes no unverified assertion
that powers of the resulting partial word extend to a common global
lift.
\end{proof}
The source and target domains in this theorem are obtained from the
generating sheets, rather than being phase charts supplied with
component labels. Their invariance and graph consistency are genuine
hypotheses expressed in observable quantities. On compact subsets
with the collars and inverse-Jacobian bounds of
Theorem~\ref{thm:v83-return}, the word map remains locally Lipschitz
under perturbation of its input lifts. A small compatibility residual
is not, by itself, asserted to imply a nearby exact common lift.

The orbit-localized ambiguity of Theorem~\ref{thm:v83-unobserved}
explains why an unsampled open orbit region cannot in general be
filled in smoothly from these graph tests. It does not show that each
particular collar, inverse-domain choice or common-lift certificate is
necessary. The positive conclusion here is instead the exact domain
recursion and a common-lift criterion that can be checked on recovered
objects.

\subsection{Selective boundary travel-time observations}
The next application uses a familiar geometric rigidity problem,
not a detector specifically engineered for the polynomial shear.
A two-arm apparatus has a physical propagation arm and a reference
arm with an unknown constant delay. Neither the readout channels,
relative throughput nor selective detector coefficients are calibrated.
The existence of the reference arm and a coarse delay bound are
apparatus information; its delay value is recovered from the records.

\begin{theorem}[Two-arm boundary-distance rigidity]
\label{thm:v84-surface}
Let $(M,g)$ be a compact simple Riemannian surface with identified
boundary. At each ordered pair $x\ne y$ of boundary points, suppose
two-arm observations have the model \eqref{eq:v83-model} with actions
\[
                     W_1(x,y)=d_g(x,y),\qquad W_2(x,y)=\tau_* ,
\]
where the constant $\tau_*$ is unknown and
$\sup_{x,y}d_g(x,y)<\tau_*<T_1$. Both arms have positive throughput,
and the two unknown conditionally independent readout channels have
full column rank. The branchwise polynomial log-detectors have known
degree bound $q$, with otherwise unknown coefficients that may depend
on $x,y$. Then $q+3$ projective deadline matrices at every such pair
determine the boundary distance, the unknown reference delay and the
metric up to a boundary-fixing isometry.
\end{theorem}
\begin{proof}
The two actions are unequal at every pair. Theorem~\ref{thm:v83-main}
therefore recovers their unordered values and both unknown channels
from the projective records. The smaller recovered value is the
boundary distance, by the declared delay ordering; this use of order
is justified by an apparatus separation, not by sorting sheets across
an action crossing. The larger value gives $\tau_*$. Continuity adds
the zero boundary distance on the diagonal. The boundary-distance
rigidity theorem for simple surfaces \cite{V84PU} then identifies the
metric up to a diffeomorphism fixing the boundary. In particular two
admissible metrics with the same projective records have identical
boundary distances, even if their detector coefficients, throughputs,
channels and initially unknown reference delays differ.
\end{proof}
To realize the clock factor, choose a release delay uniform on
$[0,D]$ with $D>T_J$, independently of the arm, and use an apparatus
bound placing both travel times in $[0,T_1)$. The event
$\alpha+W_b\le T_j$ has probability $(T_j-W_b)/D$. On a compact gate,
a positive unknown source and the endpoint change of variables give
a positive arm weight. A constant-delay reference path supplies the
second weight. Apply the branchwise retention and the two independent
readouts as in Proposition~\ref{prop:v83-physical}; a small common
baseline keeps retention below one. The same source is used at all
clocks. This construction uses no sensor response tied to a special
formula for the unknown metric.

The geometric last step is the cited boundary-distance theorem, not a
new proof of that theorem or a claim for nonsimple surfaces. The new
input reduction removes detector and channel calibration while making
the required reference acquisition explicit. This application coexists
with the nonconvex multiple-sheet example: it does not replace or
remove that action-crossing regime. For a flow, reconstructing a
suspension on a covered invariant section domain still gives only that
suspension; the whole contact manifold follows only when the whole
relevant global section and its return trajectories are covered.
